\section{\SBT dictionary for mathematical closures}
\label{sec:dictionary}

This section adapts the minimal working dictionary from the Six-Birds framework \cite{SixBirdsFramework} into a mathematical setting.
We use it as a naming convention for the rest of the paper: each ``bird'' names a role that must be handled explicitly when a discrete protocol is compressed into a stable macro-level law.

\subsection{Six primitives as closure mechanics}
We follow the framework convention that the six primitives are not ``topics'' but \emph{roles} played by concrete objects in an instantiation.
In mathematics, the relevant objects are typically families indexed by a refinement parameter (meshes, truncations, partitions), defect bounds (error terms, residual norms), quotient/completion maps, and coherence constraints (linearity, locality, invariance, compatibility with algebraic operations).

\begin{remark}[Six primitives (P1--P6), in framework language]\label{rem:six-primitives}
Fix a space $\mathcal P$ of protocol objects (discrete descriptions) equipped with a refinement chain, modeled as a family of equivalences $\sim_j$ indexed by a stage parameter $j$ (or, in analysis notation, by a step size $h\downarrow 0$).
Write $\Pi_j:\mathcal P\to \mathcal P/\!\sim_j$ for the quotient/packaging map.
Then the six primitives may be read as follows:
\begin{enumerate}
  \item \Pfive\ (Packaging). Each equivalence $\sim_j$ yields a quotient/packaging map $\Pi_j$ and induces an idempotent saturation on predicates/events via
  \[
    \mathrm{cl}_j(A)\;:=\;\Pi_j^{-1}(\Pi_j(A)).
  \]
  This is a standard closure-operator construction in order/lattice settings; see e.g.\ \cite{DaveyPriestley2002}.
  In analytic instantiations, $\Pi_j$ is a limit/identification map (``equal up to $o(1)$ as $h\to 0$''); in algebraic instantiations it is a literal quotient by a congruence.

  \item \Psix\ (Accounting). The refinement order on quotients induces a preorder and monotone defect/audit quantities (e.g.\ $j\mapsto \mathrm{Def}(j)$ or $h\mapsto \mathrm{Err}(h)$).
  In this mathematics instantiation, ``accounting'' is a feasibility ledger (error control), not a directionality certificate.

  \item \Pfour\ (Staging). The refinement chain generates the stage index $j$ (or $h$).
  Nontrivial staging corresponds to strict refinement at some scale (a genuine increase in resolution rather than a reparameterization).

  \item \Ptwo\ (Constraints). Feasible macrostates at depth $j$ are those representable as images under $\Pi_j$.
  Concretely, this is where invariances (translation/rotation/scale), locality, linearity, and algebra-compatibility are imposed as coherence conditions on the packaged layer.

  \item \Pone\ (Operator rewrite). Fix a micro-update or micro-operator $F:\mathcal P\to\mathcal P$.
  The induced macro-update $F^\sharp([p]) := [F(p)]$ on $\mathcal P/\!\sim_j$ is well-defined iff
  \[
    p\sim_j q \ \Rightarrow\ F(p)\sim_j F(q).
  \]
  When this fails, one must either refine the lens (change $\sim_j$), modify/extend $F$ (renormalize / rescale / rewrite it), or accept that no closed macro dynamics exists at depth $j$.

  \item \Pthree\ (Protocol / holonomy). When there exist two admissible closure/evolution routes with the same input/output type (e.g.\ ``refine then apply'' versus ``apply then refine''), their discrepancy defines a holonomy-like diagnostic.
  In normed instantiations we measure this by a route-mismatch quantity $\mathrm{RM}$: a number that should decay under refinement in stable regimes.
\end{enumerate}
\end{remark}

\subsection{``A mathematics is A theory'': porting the theory-package schema}
The framework packages an instantiation as a \emph{theory package}: a carrier (micro) space, a lens of observation, a completion/packaging rule, and an audit/defect principle \cite{SixBirdsFramework}.
For the mathematics instantiation we use the same shape, with ``audit'' interpreted as a nonnegative defect ledger (error and mismatch diagnostics).

\begin{definition}[Mathematics theory package]\label{def:math-theory-package}
A \emph{mathematics theory package} is a tuple
\[
\mathcal T \;=\; (Z,\ f,\ \Sigma_f,\ E,\ \mathcal D)
\]
where:
\begin{itemize}
  \item $Z$ is a carrier of discrete protocol objects at a chosen stage (e.g.\ grid functions, truncated sums/products, finite compositions, Cauchy sequences of rationals).
  \item $f:Z\to X$ is a \emph{lens} (a coarse description map). The induced definability structure $\Sigma_f$ is the family of predicates on $Z$ that are visible through $f$ (equivalently: constant on fibers of $f$).
  \item $E:\mathcal V\to\mathcal V$ is a \emph{completion/packaging endomap} on a chosen description space $\mathcal V$ associated to $Z$ (often idempotent or approximately idempotent in a specified norm). The fixed points of $E$ are the internally complete objects recognized by the theory at the chosen scale.
  \item $\mathcal D$ is a \emph{defect ledger}: a designated family of nonnegative diagnostics (error bounds, residual norms, route mismatch statistics) together with a monotonicity principle under refinement/coarsening. In this mathematics instantiation, $\mathcal D$ records feasibility (stability under refinement), not a thermodynamic arrow-of-time.
\end{itemize}
\end{definition}

\begin{remark}[Slogan]\label{rem:math-is-theory}
The title phrase ``A mathematics is A theory'' is used in the following technical sense:
to introduce a stable mathematical layer is to specify (i) what discrete protocol objects are being staged, (ii) what identifications are permitted under refinement, (iii) what completion/packaging map is being used, and (iv) what defect ledger certifies that the packaged layer is coherent.
\end{remark}

\subsection{Notation used in this instantiation}
We use $h>0$ for a refinement parameter (mesh size / step size) and write $h\downarrow 0$ for infinite refinement; equivalently we use a depth index $j\to\infty$.
Families of discrete objects are written with stage subscripts (e.g.\ $A_h$, $F_h$).
A generic route mismatch diagnostic takes the form
\[
\mathrm{RM}(h)\;=\;\frac{\|R_1(h)-R_2(h)\|}{\|R_1(h)\|+\varepsilon},
\]
where $R_1,R_2$ are two admissible routes and $\varepsilon>0$ is a small stabilizer.
In later exhibits, $\mathrm{RM}$ is instantiated as: (i) a coordinate-change mismatch for discrete derivatives, and (ii) a mismatch between two staged prime closures after prototype packaging.
